% GOVERNANCE NOTE::LSPV5-WRITE-REPAIR::2026-10-12T08:17:14Z
\documentclass[11pt]{article}
\usepackage[margin=1in]{geometry}
\usepackage{xurl,hyperref}
\begin{document}
\section*{Governance note: LSPV-5 canonical-write repair}
\textbf{Change ID:} LSPV5-WRITE-REPAIR-20261012-0817.

\textbf{Evidence.} GitHub preflight found
\path{research/main/2026-10-12_0448_MAIN_SYNTHESIS-0022.tex} as the newest
canonical LSPV-5 artifact.  It was substantive and had the required MAIN
header and final sentinel.  The enabled LSPV-5 automation showed later
completed invocations but no newer canonical report existed, spanning at
least two scheduled cycles.

\textbf{Classification.} \textbf{PROMPT / PROCEDURE CONTRADICTION.}  The
prompt requested a timestamped SHA-safe GitHub/Drive fallback but omitted the
canonical \path{research/main/YYYY-MM-DD_HHMM_MAIN_<RUN_ID>.tex} path,
mandatory header/sentinel readback, and same-run current-SHA reconciliation of
root \path{main.tex} required by \path{AGENTS.md}.

\textbf{Repair.} The automation schedule, enabled state, title and scientific
target were preserved.  Its prompt received one appended predecessor-recovery
paragraph requiring canonical-path retry of an available exact body, GitHub
readback, required markers and current-SHA root reconciliation.  Missing
scientific content may not be reconstructed.

\textbf{Rollback.} Remove only the appended paragraph beginning
``Predecessor write recovery'' from automation
\texttt{6ac44a4166a88191b3b582d8909b11da}.  No schedule or scientific text
needs rollback.

\textbf{Status.} Repair applied and automation left enabled.  This note does
not assert that a missing scientific body existed or was recovered.

% END_OF_GOVERNANCE_NOTE::LSPV5-WRITE-REPAIR-20261012-0817
\end{document}
